\subsection{対称性と漸化式}

$u=1-x$と置くと2つの指数が入れ替わるため
\begin{equation}\label{式:ベータ対称性}
 B(p,q)=B(q,p)
\end{equation}
である.広義積分の場合も,端点を避けた区間で置換し,各端点の極限を取れば同じ結論になる.

次に,整数の場合と同じく,指数を1つずらした積分どうしの関係を作ろう.被積分関数に$x+(1-x)=1$を掛ければ
\begin{equation}\label{式:ベータ和の関係}
 B(p+1,q)+B(p,q+1)=B(p,q)
\end{equation}
となる.もう1つの関係を得るために$x^p(1-x)^q$を微分する.
\[
 \frac{\dd}{\dd x}\{x^p(1-x)^q\}
 =p x^{p-1}(1-x)^q-qx^p(1-x)^{q-1}
\]
これを$\varepsilon$から$1-\delta$まで積分し,両端をそれぞれ極限へ近づける.\,$p,q>0$なので$x^p(1-x)^q$は両端で$0$へ近づき,右辺の2つの積分はすでに収束すると分かっている.よって
\[
 0=pB(p,q+1)-qB(p+1,q)
\]
を得る.これは整数の場合の部分積分を,端点を避けて行ったものである.式\eqref{式:ベータ和の関係}に$B(p,q+1)=\frac qpB(p+1,q)$を代入すると
\[
 \frac{p+q}{p}B(p+1,q)=B(p,q)
\]
となるので
\begin{equation}\label{式:ベータ漸化式}
 B(p+1,q)=\frac{p}{p+q}B(p,q),\qquad
 B(p,q+1)=\frac{q}{p+q}B(p,q)
\end{equation}
を得る.この漸化式は整数以外にも成り立つ.

整数のときには,\,$B(1,1)=\int_0^1 1\dd x=1$から出発して階乗の公式を再び導ける.右の式を$q=1,2,\ldots,n$について順に使うと
\[
 B(1,n+1)=\frac{n}{n+1}\cdot\frac{n-1}{n}\cdots\frac12B(1,1)=\frac{1}{n+1}
\]
となる.次に左の式を$p=1,2,\ldots,m$について順に使うと
\begin{align*}
 B(m+1,n+1)
 &=\frac{m}{m+n+1}\cdot\frac{m-1}{m+n}\cdots\frac{1}{n+2}B(1,n+1)\\
 &=\frac{m!\,n!}{(m+n+1)!}
\end{align*}
となり,式\eqref{式:整数ベータ積分}と一致する.どちらも比が添字の式で与えられる階比型であり,比を順に掛けて一般項を得ている.
